\documentclass[../../../include/open-logic-section]{subfiles}

\begin{document}

\olfileid[fr]{sth}{choice}{woproblem}
\olsection{Le problème du bon ordre}

Il est clair que beaucoup de choses dépendent
de notre acceptation de l'axiome du bon ordre.
Mais la discussion s'est surtout concentrée
sur un principe équivalent : l'axiome du choix.
Nous allons donc nous tourner vers celui-ci
et démontrer cette équivalence.

En \citeyear{Cantor1883}, Cantor exprimait
son soutien à l'axiome du bon ordre, qu'il
appelait «~une loi de la pensée qui me
paraît fondamentale, riche de conséquences
et particulièrement remarquable par sa
validité générale~» (cité par
\citeauthor{Potter2004}
\citeyear[p.~243]{Potter2004}).
Mais Cantor finit par se convaincre
que cet «~axiome~» devait être démontré.
La communauté mathématique en vint
à la même conclusion.

Le problème fut «~résolu~» par Zermelo
en \citeyear{Zermelo1904}. Pour expliquer
sa solution, il nous faut quelques définitions.

\begin{defn}
Une fonction $f$ est une \emph{fonction de choix}
si et seulement si $f(x) \in x$ pour tout
$x \in \dom{f}$. Nous disons que $f$ est
une \emph{fonction de choix pour $A$} si
$f$ est une fonction de choix telle que
$\dom{f} = A \setminus \{\emptyset\}$.
\end{defn}

Intuitivement, pour chaque ensemble non vide
$x \in A$, une fonction de choix pour $A$
\emph{choisit} un élément particulier
$f(x)$ de $x$. Voici donc l'axiome du choix :

\begin{axiom}[Choix]
	Tout ensemble admet une fonction de choix.
\end{axiom}

Zermelo a montré que le choix entraîne
le bon ordre, et réciproquement :

\begin{thm}[dans $\ZF$]\ollabel{thmwochoice}
L'axiome du bon ordre et l'axiome du choix
sont équivalents.
\end{thm}

\begin{proof}
\emph{De gauche à droite.}
Soit $A$ un ensemble d'ensembles.
L'ensemble $\bigcup A$ existe par l'axiome
de l'union ; d'après l'axiome du bon ordre,
il existe donc un ordre $<$ qui
bien ordonne $\bigcup A$. Posons
$f(x) = \text{le plus petit élément de $x$
pour l'ordre $<$}$.
C'est une fonction de choix pour $A$.

\emph{De droite à gauche.}
Fixons $A$. Si $A=\emptyset$, il est déjà
bien ordonné. Supposons donc $A$ non vide.
Par l'axiome du choix, il existe une fonction
de choix $f$ pour
$\Pow{A} \setminus \{\emptyset\}$.
Le théorème de récursion donne d'abord un \emph{terme} pour les
valeurs successives (voir les remarques qui suivent
le \olref[sth][spine][recursion]{transrecursionschema}).
D'après le \olref[hartogs]{HartogsLemma}, il existe un ordinal
$\theta$ qui ne s'injecte pas dans $A$. Fixons un marqueur
$\ast\notin A$. Sur l'ordinal $\theta$, le remplacement permet de
rassembler les valeurs de la récursion en une fonction $g$ de
domaine $\theta$, définie pour tout $\beta\in\theta$ par
\[
g(\beta) =
  \begin{cases}
    f(A\setminus\funimage{g}{\beta})
      & \text{si } A\setminus\funimage{g}{\beta}\neq\emptyset,\\
    \ast & \text{sinon}.
  \end{cases}
\]

Avant la première occurrence du marqueur, chaque valeur choisie
appartient à $A$ et diffère de toutes les valeurs précédentes.
Si le marqueur n'apparaissait jamais, $g$ serait une injection de
$\theta$ dans $A$, en contradiction avec le choix de $\theta$.

Soit donc $\alpha$ le premier ordinal pour lequel $g(\alpha)=\ast$.
Les valeurs précédentes sont toutes dans $A$ et les éléments de $A$
sont tous atteints, puisque $A\setminus\funimage{g}{\alpha}=\emptyset$.
La restriction $g\!\upharpoonright\!\alpha$ est ainsi une
!!{bijection} de l'ordinal $\alpha$ sur $A$. Le bon ordre de
$\alpha$ se transporte par cette bijection sur $A$.
\end{proof}

Ainsi, le bon ordre et le choix
se soutiennent ou tombent ensemble.
La question demeure : devons-nous
les accepter ?

\end{document}
